\documentclass[11pt]{article}

\usepackage[margin=1in]{geometry}
\usepackage{amsmath,amssymb}
\usepackage{booktabs}
\usepackage{tabularx}
\usepackage{array}
\usepackage{enumitem}
\usepackage[T1]{fontenc}
\usepackage{hyperref}

\newcommand{\code}[1]{\texttt{#1}}
\emergencystretch=3em
\newcolumntype{L}{>{\raggedright\arraybackslash}X}

\title{\code{cluster\_diagnostics.py}: cluster statistics and radial\\
distribution functions for 2D AMReX plotfiles}
\author{FHDeX}
\date{}

\begin{document}
\maketitle

\section{Overview}

\code{cluster\_diagnostics.py} (in \code{FHDeX/python}) post-processes a sequence of
2D AMReX plotfiles, written for the Dean--Kawasaki runs in
\code{exec/dean\_kow/interaction} (for example \code{inputs\_fv\_pp\_merge} and
\code{inputs\_fv\_hk\_cluster}). For each plotfile it
\begin{enumerate}
  \item finds the clusters in the density field and measures each one's mass,
        centre, size and peak density;
  \item matches them to the clusters in the previous plotfile, giving each
        cluster a stable id and logging mergers, splits, new clusters and
        clusters that dissolve;
  \item accumulates two radial distribution functions: $g(r)$ of the density
        field and $g(r)$ of the cluster centres, and the structure factor $S(k)$;
  \item measures the orientational order of the cluster centres ($\psi_6$)
        and whether $S(k)$ has Bragg spots or a ring.
\end{enumerate}
The results are whitespace-separated text tables with \code{\#} header lines.
The script needs only Python~3 and numpy. It reuses the plotfile reader in
\code{stack\_plotfiles\_time.py}, which must be in the same directory.

\paragraph{Usage.}
\begin{verbatim}
python3 cluster_diagnostics.py plt_pp_* --num-part 327680 -o diag
python3 cluster_diagnostics.py plt_pp_* --num-part 327680 -o diag \
        --t-range 0.02 0.03 --rmax 0.4
\end{verbatim}
Plotfiles are sorted by the step number in their headers, not by file name.
All of them must have the same domain and cell size.

\section{Assumptions and limitations}

\begin{itemize}
  \item \textbf{2D, periodic domain.} Clusters may wrap across the boundaries,
        and all distances use the minimum-image convention. The plotfile
        header does not record periodicity, so the script assumes it.
  \item \textbf{Level 0 only.} Finer levels are ignored, with a warning.
  \item \textbf{Density field only.} The plotfiles contain no particle
        positions, so both $g(r)$ are computed from the density variable
        (\code{phi0} by default).
  \item \textbf{Touching clusters count as one.} Two clusters become one as soon
        as the density between them exceeds the threshold. This is the
        definition of a merger used here; clusters are not split at density
        saddles.
  \item \textbf{Total mass includes negative densities.} At low particle counts
        per cell $\phi$ can be slightly negative between clusters. The total
        mass $\sum\phi\,\Delta V$, used for \code{-{}-mmin} and the mass fraction,
        includes those values, so the mass fraction in clusters can slightly
        exceed 1.
\end{itemize}

\section{Definitions}

Let the domain be $[x_{\mathrm{lo}}, x_{\mathrm{lo}}+L_x)\times[y_{\mathrm{lo}}, y_{\mathrm{lo}}+L_y)$
with cells of size $\Delta x\times\Delta y$, area $\Delta V = \Delta x\,\Delta y$
and centres $\mathbf{x}_i$, and let $\phi_i$ be the density in cell $i$. The
minimum-image displacement in direction $d$ is
\begin{equation}
  \operatorname{mi}_d(a) = \Bigl(a + \tfrac12 L_d\Bigr) \bmod L_d \;-\; \tfrac12 L_d
  \in \bigl[-\tfrac12 L_d,\ \tfrac12 L_d\bigr).
\end{equation}

\subsection{Cluster}

With the mean density $\bar\phi = \sum_i \phi_i\,\Delta V/(L_xL_y)$ and the
threshold factor $c$ (\code{-{}-threshold}), a cluster is a set of cells connected
through shared faces, including across the periodic boundaries, in which
\begin{equation}
  \phi_i > c\,\bar\phi .
\end{equation}
It is kept only if its mass $M = \sum_{i\in C}\phi_i\,\Delta V$ is at least
$m_{\min}$ (\code{-{}-mmin}) times the total mass. For each cluster $C$, with
weights $w_i = \phi_i\,\Delta V$:

\begin{center}
\begin{tabularx}{\textwidth}{@{}lL@{}}
\toprule
Quantity & Definition \\
\midrule
mass $M$ & $\sum_{i\in C} w_i$ \\
centre $x_c$ (and $y_c$) &
  Circular mean in each direction,
  $x_c = \dfrac{L_x}{2\pi}\operatorname{atan2}\!\Bigl(\sum w_i\sin\tfrac{2\pi x_i}{L_x},\
  \sum w_i\cos\tfrac{2\pi x_i}{L_x}\Bigr) \bmod L_x$. Correct for clusters that
  straddle the boundary. \\
$R_g$ & Radius of gyration,
  $R_g^2 = \sum w_i\,|\mathbf{d}_i|^2/M$ with
  $\mathbf{d}_i = \bigl(\operatorname{mi}_x(x_i - x_c),\ \operatorname{mi}_y(y_i - y_c)\bigr)$ \\
peak & $\max_{i\in C}\phi_i$ \\
ncells & number of cells in $C$ \\
wide & 1 if some cell has $|d_{i,x}| \ge 0.45\,L_x$ or $|d_{i,y}| \ge 0.45\,L_y$.
  The minimum image, and so the centre and $R_g$, are then ambiguous; this
  mainly happens for a single cluster that fills the box. \\
\bottomrule
\end{tabularx}
\end{center}

\subsection{Tracking between frames}

For consecutive plotfiles with old clusters $b$ and new clusters $a$, let
$O_{ba} = \sum_{i\in b\cap a}\phi^{\mathrm{old}}_i\,\Delta V$ be the part of
the old cluster's mass, at the old time, in cells that now belong to $a$.
\begin{enumerate}
  \item \textbf{Matching.} Old cluster $b$ is assigned to
        $\arg\max_a O_{ba}$ if $\max_a O_{ba} > 0.1\,M_b$. Otherwise it is
        assigned to the new cluster with the nearest centre (minimum image),
        provided that distance is at most \code{-{}-search-radius}. The
        fallback catches the end of a merger, when a small cluster falls into
        a large one within one frame and its old cells no longer overlap
        anything.
  \item \textbf{Merger.} A new cluster $a$ with two or more old clusters
        assigned is a merger, provided $M_a \ge 0.7\sum_{b\to a} M_b$.
        Assigned old clusters that would break this condition are treated as
        dissolved instead, starting with the farthest.
  \item \textbf{Split.} An old cluster with $O_{ba} > 0.1\,M_b$ for two or more
        new clusters $a$ is a split.
  \item \textbf{Formed and dissolved.} A new cluster with no old cluster
        assigned, and not split off one, is formed. An old cluster assigned to
        nothing, and not split, is dissolved.
  \item \textbf{Ids.} A new cluster takes the id of its most massive assigned
        old cluster. Formed clusters and pieces split off get new ids, so ids
        are never reused.
\end{enumerate}
The thresholds 0.1 and 0.7 are fixed in the code
(\code{track(frac=0.1, mass\_keep=0.7)}). Tracking needs plotfiles close enough
in time that clusters move less than about their own size between frames,
apart from the final plunge of a merger.

\subsection{\texorpdfstring{Density-field $g(r)$}{Density-field g(r)}}

From the cell particle counts $n_i = \phi_i\,N\,\Delta V$, with
$N$ = \code{-{}-num-part}, $M$ the number of cells and $N_{\mathrm{tot}} = \sum_i n_i$:
\begin{align}
  G(\boldsymbol{\Delta}) &= \sum_i n_i\,n_{i+\boldsymbol{\Delta}}
     - \delta_{\boldsymbol{\Delta},0}\sum_i n_i , \\
  g(\boldsymbol{\Delta}) &= \frac{M\,G(\boldsymbol{\Delta})}{N_{\mathrm{tot}}(N_{\mathrm{tot}}-1)} .
\end{align}
$G$ is the periodic autocorrelation, computed by FFT, with the self pairs
removed at $\boldsymbol{\Delta}=0$. A Poisson (ideal gas) field gives
$g = 1$ at every $\boldsymbol{\Delta}$, including $\boldsymbol{\Delta}=0$. Each
lattice displacement $\boldsymbol{\Delta}$ is wrapped componentwise into
$[-M_d/2, M_d/2)$, so $r = |(\Delta_x\Delta x,\ \Delta_y\Delta y)|$ is the
minimum-image distance. The value of $g(r)$ in a bin $[r_k, r_k+\Delta r)$ is
the average of $g(\boldsymbol{\Delta})$ over the displacement vectors that fall
in it. This normalizes by the exact number of lattice vectors per shell rather
than by $2\pi r\,\Delta r$. The result is averaged over the frames in
\code{-{}-t-range}. $N$ enters only through the self-pair correction, which has
relative size $1/N$.

\subsection{\texorpdfstring{Cluster-centre $g(r)$}{Cluster-centre g(r)}}

For every frame in \code{-{}-t-range} with $n_c\ge 2$ clusters, the minimum-image
distances between all pairs of cluster centres are histogrammed, counting each
pair once. With domain area $A = L_xL_y$,
\begin{equation}
  g(r_k) = \frac{\sum_{\text{frames}} \#\{\text{pairs with } r\in[r_k, r_k+\Delta r)\}}
                {\sum_{\text{frames}} \tfrac12 n_c(n_c-1)\,\pi\bigl((r_k+\Delta r)^2 - r_k^2\bigr)/A} .
\end{equation}
The annulus area is exact for $r \le L/2$, because every such annulus lies
inside the minimum-image cell. Pair counts and expected counts are summed over
frames before dividing, because single frames often hold only about ten
clusters.

\subsection{\texorpdfstring{Structure factor $S(k)$}{Structure factor S(k)}}
From the same cell counts as the field $g(r)$,
\begin{equation}
  S(\mathbf k) = \frac{|\hat n(\mathbf k)|^2}{N_{\mathrm{tot}}}, \qquad
  \hat n(\mathbf k) = \sum_i n_i\,e^{-i\mathbf k\cdot\mathbf x_i}, \qquad \mathbf k\ne0 ,
\end{equation}
so that a Poisson field gives $S=1$ and
$S(k)=1+(N/V)\int(g-1)e^{i\mathbf k\cdot\mathbf r}d\mathbf r$. Wave vectors come
from \code{fftfreq}: index $i>N_d/2$ is the negative wavenumber $i-N_d$, so
$k_d=2\pi\cdot(\text{index})/L_d$. $S$ is averaged over shells
$|\mathbf k|\in[(j-\tfrac12)\Delta k,(j+\tfrac12)\Delta k)$ with
$\Delta k$ = \code{-{}-dk} (default $2\pi/\max L_d$), each divided by its exact
number of modes. Only complete shells up to the smallest Nyquist wavenumber
$\pi/\Delta x_d$ are written, averaged over the frames in \code{-{}-t-range}.
For many samples use the in-code structure factor of
\code{src\_deankow/interaction} (\code{struct\_fact\_int}), which writes the same
shell average.

\subsection{Orientational order: crystal or amorphous}
$g(r)$ averages over directions, so a cluster crystal and an amorphous packing
with the same spacing can have the same $g(r)$ peaks. Two measures, written to
\code{order.txt} for every frame, tell them apart.

\paragraph{Hexatic bond order.} Each cluster centre $i$ is joined to its
\code{-{}-nn} nearest neighbours (default 6, minimum image), or to all
neighbours closer than \code{-{}-nn-cut}. With bond angles $\theta_{ij}$,
\begin{equation}
  \psi_{6,i} = \frac1{n_i}\sum_j e^{6i\theta_{ij}} ,\qquad
  \psi_6^{\mathrm{loc}} = \overline{|\psi_{6,i}|},\qquad
  \psi_6^{\mathrm{glob}} = \bigl|\,\overline{\psi_{6,i}}\,\bigr| .
\end{equation}
A hexagonal crystal gives 1 for both. Random points give
$\psi_6^{\mathrm{loc}}\approx0.36$ and $\psi_6^{\mathrm{glob}}\approx0.4/\sqrt N$.
A high local but low global value means hexagonal grains with different
orientations.

\paragraph{Bragg spots or a ring.} On the shell with the largest mean $S$ of a
frame's full $S(\mathbf k)$, with $n$ modes,
\begin{equation}
  \mathrm{PR} = \frac{\bigl(\sum S\bigr)^2}{n\sum S^2} ,
\end{equation}
which is $\approx(\text{number of spots})/n$ for Bragg spots (6/60 = 0.100 for a
hexagonal lattice of blobs) and $\approx0.5$ for an isotropic ring with random
phases. \code{top\_share} is the fraction of the shell's $S$ in its
\code{-{}-top-modes} largest modes.

\section{Options}

\begin{center}
\begin{tabularx}{\textwidth}{@{}l l L@{}}
\toprule
Option & Default & Meaning \\
\midrule
\code{plotfiles} & (required) & 2D plotfile directories, e.g.\ \code{plt\_pp\_*} \\
\code{-{}-num-part N} & (required) & \code{num\_part} of the run; converts $\phi$ to particle counts for the field $g(r)$ \\
\code{-{}-var NAME} & \code{phi0} & plotfile variable used as the density \\
\code{-{}-threshold C} & 5 & cluster cells have $\phi > C\,\bar\phi$ \\
\code{-{}-mmin F} & 0.01 & smallest cluster kept, as a fraction of the total mass \\
\code{-{}-every K} & 1 & use every $K$-th plotfile, after sorting by step \\
\code{-{}-rmax R} & $L/2$ & range of both $g(r)$; values above $L/2$ (shortest side) are reduced to $L/2$ with a warning \\
\code{-{}-dr D} & $\Delta x$ & bin width of the field $g(r)$; with $D\le\Delta x$ the first bin holds only $\boldsymbol{\Delta}=0$ \\
\code{-{}-dr-clusters D} & $L/50$ & bin width of the cluster-centre $g(r)$ \\
\code{-{}-search-radius R} & $L/4$ & largest distance a centre may move between frames and still be matched when its old cells overlap no new cluster \\
\code{-{}-t-range T0 T1} & all frames & time window for both $g(r)$ averages; cluster statistics and tracking always use every frame \\
\code{-{}-rdf-per-frame} & off & also write the field $g(r)$ of each frame in the window \\
\code{-{}-dk D} & $2\pi/\max L$ & shell width of the structure factor \\
\code{-{}-sf-per-frame} & off & also write the per-frame shell-averaged $S(k)$ \\
\code{-{}-nn K} & 6 & bond order: number of nearest neighbours \\
\code{-{}-nn-cut R} & off & bond order: all neighbours closer than $R$ instead \\
\code{-{}-top-modes K} & 6 & \code{order.txt} \code{top\_share}: largest modes on the $S(k)$ peak shell \\
\code{-{}-sf-full} & off & write the full $S(\mathbf k)$, averaged over the window, as the plotfile \code{sf\_full\_avg} \\
\code{-{}-sf-full-per-frame} & off & also write the full $S(\mathbf k)$ of each frame as \code{sf\_full\_<step>} \\
\code{-o DIR}, \code{-{}-output DIR} & \code{cluster\_diag} & output directory, created if needed; existing files are overwritten \\
\bottomrule
\end{tabularx}
\end{center}
The script also prints one line per frame with the cluster count, the mass
fraction in clusters and the events since the previous frame.

\section{Output files}

Times are plotfile times; lengths and positions are in the plotfile's physical
units. Masses are $\sum\phi\,\Delta V$, a fraction of the total, because the
codes normalize $\phi$ to integrate to 1.

\subsection*{\code{summary.txt}: one line per plotfile}
\begin{center}
\begin{tabularx}{\textwidth}{@{}l L@{}}
\toprule
Column & Meaning \\
\midrule
\code{t}, \code{step} & time and step \\
\code{nclusters} & number of clusters \\
\code{mass\_fraction} & total mass in clusters / total mass \\
\code{largest\_mass} & mass of the largest cluster (0 if none) \\
\code{merges}, \code{splits} & merger and split events since the previous plotfile \\
\code{forms}, \code{dissolves} & clusters formed and dissolved since the previous plotfile \\
\bottomrule
\end{tabularx}
\end{center}

\subsection*{\code{clusters.txt}: one line per cluster per plotfile, sorted by id}
\begin{center}
\begin{tabularx}{\textwidth}{@{}l L@{}}
\toprule
Column & Meaning \\
\midrule
\code{t}, \code{step} & time and step \\
\code{id} & cluster id, stable across frames \\
\code{mass} & cluster mass $M$ \\
\code{xc}, \code{yc} & centre of mass, in $[\text{prob\_lo}, \text{prob\_hi})$ \\
\code{Rg} & radius of gyration \\
\code{peak} & largest $\phi$ in the cluster \\
\code{ncells} & number of cells \\
\code{wide} & 1 if the cluster spans about half the box or more ($x_c$, $y_c$, $R_g$ unreliable) \\
\code{psi6} & $|\psi_{6,i}|$ of the cluster (NaN without neighbours) \\
\bottomrule
\end{tabularx}
\end{center}
To follow one cluster, filter by id, e.g.\ \code{awk '\$3==4' clusters.txt}.

\subsection*{\code{events.txt}: one line per merger or split}
\begin{center}
\begin{tabularx}{\textwidth}{@{}l L@{}}
\toprule
Column & Meaning \\
\midrule
\code{t}, \code{step} & time and step of the first plotfile in which the event shows \\
\code{type} & \code{merge} or \code{split} \\
\code{old\_ids} & comma-separated ids before the event: the merging clusters, or the cluster that split \\
\code{old\_masses} & their masses at the previous plotfile \\
\code{new\_ids} & id(s) after the event: the merged cluster, or the pieces of the split \\
\code{new\_masses} & their masses \\
\bottomrule
\end{tabularx}
\end{center}
Formed and dissolved clusters appear only as counts in \code{summary.txt}.

\subsection*{\texorpdfstring{\code{rdf\_field.txt}: density-field $g(r)$}{rdf_field.txt}}
The first header line gives the number of frames, the time window,
\code{num\_part}, $\Delta r$ and $r_{\max}$. Then one line per non-empty bin:
\begin{center}
\begin{tabularx}{\textwidth}{@{}l L@{}}
\toprule
Column & Meaning \\
\midrule
\code{r} & mean distance of the lattice displacements in the bin \\
\code{g} & $g(r)$ averaged over the frames in the window \\
\code{ndisplacements} & number of lattice displacements in the bin \\
\code{g(t=...)} & with \code{-{}-rdf-per-frame}, one column per frame \\
\bottomrule
\end{tabularx}
\end{center}
The first bin, at $r=0$, is the same-cell value; in clustered states it is
large because of the cluster's own density.

\subsection*{\texorpdfstring{\code{rdf\_clusters.txt}: cluster-centre $g(r)$}{rdf_clusters.txt}}
The first header line gives the number of frames, the time window, $\Delta r$
and $r_{\max}$. Then one line per bin:
\begin{center}
\begin{tabularx}{\textwidth}{@{}l L@{}}
\toprule
Column & Meaning \\
\midrule
\code{r\_lo}, \code{r\_hi} & bin edges \\
\code{g} & $\sum$ pairs $/ \sum$ expected (0 where nothing is expected) \\
\code{pairs} & cluster pairs in the bin, summed over frames \\
\code{expected} & ideal-gas expectation for the bin, summed over frames \\
\bottomrule
\end{tabularx}
\end{center}
\code{pairs} shows how much each bin's $g$ rests on; with a few clusters per
frame, bins with only a handful of pairs are noisy.

\subsection*{\texorpdfstring{\code{sf.txt}: shell-averaged structure factor}{sf.txt}}
A header with the number of frames, the time window, \code{num\_part} and
$\Delta k$, then one line per shell: \code{k} (mean $|\mathbf k|$ of the
shell), \code{S}, \code{nmodes}, and with \code{-{}-sf-per-frame} one
\code{S(t=...)} column per frame.

\subsection*{\texorpdfstring{\code{order.txt}: orientational order, one line per plotfile}{order.txt}}
A header with the neighbour rule and \code{-{}-top-modes}, then one row per
plotfile:
\begin{center}
\begin{tabularx}{\textwidth}{@{}r l L@{}}
\toprule
\# & Column & Meaning \\
\midrule
1 & \code{t} & plotfile time \\
2 & \code{step} & plotfile step \\
3 & \code{nclusters} & number of clusters in the frame (as in \code{summary.txt}) \\
4 & \code{psi6\_local} & $\psi_6^{\mathrm{loc}}=\overline{|\psi_{6,i}|}$: how hexagonal each neighbourhood is, whatever its orientation; 1 for a perfect hexagon, $\approx0.36$ for random points \\
5 & \code{psi6\_global} & $\psi_6^{\mathrm{glob}}=|\overline{\psi_{6,i}}|$: one orientation across the box; $\approx\psi_6^{\mathrm{loc}}$ for a single crystal, $\approx0.4/\sqrt N$ with no common orientation \\
6 & \code{kpeak} & mean $|\mathbf k|$ of the $S(k)$ ring with the largest mean $S$ in this frame, normally the main peak at $\approx2\pi/\text{spacing}$ \\
7 & \code{nmodes} & number of FFT wave vectors $n$ on that ring \\
8 & \code{PR} & $(\sum S)^2/(n\sum S^2)$, roughly the fraction of modes carrying the signal: $\approx$ spots$/n$ for Bragg spots, $\approx0.5$ for a uniform ring \\
9 & \code{top\_share} & fraction of the ring's $S$ in its \code{-{}-top-modes} largest modes (default 6): $\approx1$ for a crystal, small for a ring \\
\bottomrule
\end{tabularx}
\end{center}
Example, the last frame of \code{\_new\_repul\_2}:
\begin{verbatim}
0.08773803711 400000 53 0.987506 0.986861 50.306732 48 0.125774 0.996647
\end{verbatim}
At $t=0.0877$ (step 400000) there are 53 clusters. Local and global $\psi_6$
are both 0.99: one hexagonal crystal with a single orientation. The $S(k)$
peak ring at $|\mathbf k|\approx50.3$ has 48 modes, and $\mathrm{PR}=0.126\approx6/48$,
with the 6 largest modes holding 99.7\% of it: six Bragg spots. The $t=0$ row,
\code{0 0 0 nan nan 7.58 8 nan nan}, is a uniform field: the order values are
NaN, and \code{kpeak} and \code{nmodes} only describe the first ring.

\subsection*{\texorpdfstring{\code{sf\_full\_avg}, \code{sf\_full\_<step>}: full $S(\mathbf k)$ plotfiles}{sf_full plotfiles}}
With \code{-{}-sf-full} (average over the window) or \code{-{}-sf-full-per-frame}
(one per frame), the full structure factor is written as a single-box AMReX
plotfile on the $k$ grid, in the layout of the in-code \code{plt\_SF\_mag}.
\begin{itemize}[nosep]
  \item \textbf{Layout:} cell $i$ in direction $d$ is centred at
        $k_d=2\pi(i-n_d/2)/L_d$, so $k=0$ is the centre cell and FFT indices
        above $n/2$ are negative wavenumbers. The cell size is $2\pi/L_d$.
  \item \textbf{Components:} \code{S} and \code{log10S}. Bragg spots can be
        $10^4$ times the background, so \code{log10S} is easier to view.
        At $k=0$ both are set to 0.
  \item \textbf{Viewing:} VisIt, ParaView or amrvis. A crystal shows bright
        spots (six on the first ring for a 2D hexagonal lattice); an
        amorphous state shows a uniform ring. In 3D, use slices through
        $k=0$ or an isosurface of \code{log10S}.
  \item \textbf{Size:} 16 bytes per cell (33~MB at $128^3$).
\end{itemize}
Shell-averaging \code{sf\_full\_avg} reproduces \code{sf.txt}.

\section{Example}

The \code{inputs\_fv\_pp\_merge} run ($256^2$, \code{num\_part} $=327680$,
seed 11) gives:
\begin{itemize}
  \item clusters from $t\approx 0.014$, ten of them by $t\approx 0.018$;
  \item nine mergers between $t\approx 0.027$ and $0.036$, ending in one
        cluster with 96--99\% of the mass;
  \item with \code{-{}-t-range 0.02 0.03}, a cluster-centre $g(r)$ with almost
        no pairs below $r\approx 0.2$ and a peak over $0.24$--$0.36$ (expected
        spacing about $0.32$);
  \item a field $g(r)$ of about 14 at $r=0$, a dip to about 0.2 near
        $r\approx 0.13$, and a second peak of about 1.5 at $r\approx 0.3$.
\end{itemize}
The cluster tracks can be viewed in space--time by stacking the same plotfiles
with \code{stack\_\allowbreak plotfiles\_\allowbreak time.py}.

\end{document}
